% Fichier généré automatiquement par tools/generate_corrections.py.
% Source : TEC/TEC-04-Meq-Jeq/27_TrainEpi/27_TrainEpi.tex
% Statut : complété (3/3 question(s)).
\begin{corrigebox}[Corrigé complété]
\CorrigeQuestion{1}
On construit un sommet par solide, puis une arête par liaison en
        indiquant son type, son centre et son axe. Les actions extérieures (poids,
        actionneur, charge) sont reliées au solide concerné. Le graphe doit retrouver le
        nombre de mobilités et de boucles du mécanisme de l'énoncé.
\CorrigeQuestion{2}
En bloquant le porte satellite, on a : $\dfrac{\omega_{03}}{\omega_{13}}=-\dfrac{Z_1}{Z_0}$. On a donc, 
$\dfrac{\omega_{03}}{\omega_{10}+\omega_{03}}=-\dfrac{Z_1}{Z_0}$

$\Leftrightarrow \dfrac{\omega_{30}}{\omega_{30}-\omega_{10}}=-\dfrac{Z_1}{Z_0}$
$\Leftrightarrow \omega_{30}=-\dfrac{Z_1}{Z_0} \omega_{30}+\dfrac{Z_1}{Z_0}\omega_{10} $
$\Leftrightarrow \omega_{30}\left( 1+\dfrac{Z_1}{Z_0} \right)=\dfrac{Z_1}{Z_0}\omega_{10} $
$\Leftrightarrow \omega_{30}=\dfrac{Z_1}{Z_0+Z_1}\omega_{10} $.
\CorrigeQuestion{3}
On choisit les inconnues et les conventions positives de la figure,
        on écrit les relations de conservation/fermeture adaptées, puis on résout le
        système obtenu. Le résultat symbolique doit être homogène et vérifié dans les cas
        limites (entrée nulle, rapport unitaire ou position de référence).
\end{corrigebox}
